%!TEX root = ../dokumentation.tex

\addchap{\langabkverz}
%nur verwendete Akronyme werden letztlich im Abkürzungsverzeichnis des Dokuments angezeigt
%Verwendung: 
%		\ac{Abk.}   --> fügt die Abkürzung ein, beim ersten Aufruf wird zusätzlich automatisch die ausgeschriebene Version davor eingefügt bzw. in einer Fußnote (hierfür muss in header.tex \usepackage[printonlyused,footnote]{acronym} stehen) dargestellt
%		\acs{Abk.}   -->  fügt die Abkürzung ein
%		\acf{Abk.}   --> fügt die Abkürzung UND die Erklärung ein
%		\acl{Abk.}   --> fügt nur die Erklärung ein
%		\acp{Abk.}  --> gibt Plural aus (angefügtes 's'); das zusätzliche 'p' funktioniert auch bei obigen Befehlen
%       \aclp{Abk.} --> gibt die Erklärung in Plura (angefügtes 's' aus
%	siehe auch: http://golatex.de/wiki/%5Cacronym
%	
\begin{acronym}[YTMMM]
\setlength{\itemsep}{.5\parsep}

\acro{API}{App Programming Interface}
\acro{DHBW}{Duale Hochschule Baden-Würtemberg}
\acro{HVAC}{Heating, Ventilation and Air Conditioning}
\acro{MCU}{Mikrocontroller-Einheit}
\acro{EEPROM}{electrically erasable programmable read-only memory}
\acro{IoT}{Internet of Things}
\acro{OSI}{Open Systems Interconnection}
\acro{TCP}{Transmission Control Protocol}
\acro{UDP}{User Datagram Protocol}
\acro{IP}{Internet Protocol}
\acro{MAC}{Media Access Control}
\acro{ACK}{Acknowledgement}
\acro{CSMA/CD}{Carrier Sense Multiple Access with Collision Detection}
\acro{CSMA/CA}{Carrier Sense Multiple Access with Collision Avoidance}
\acro{MQTT}{Message Queuing Telemetry Transport}
\acro{HTTP}{Hypertext Transfer Protocol}
\acro{HTTPS}{Hypertext Transfer Protocol Secure}
\acro{DNS}{Domain Name System}
\end{acronym}
